\documentclass[12pt]{article} 

\usepackage[top=1.5cm, bottom=1.5cm, left=1.5cm, right=1.5cm]{geometry}
\usepackage{color,graphicx,amsmath,mathtools,amssymb,url}
\pagestyle{plain}

\begin{document}  
\pagestyle{empty}
 
Homework 3: Nico Boyle, Owen Moorhus, Asher Olgeirson




\vskip0.5in


%\noindent {\small This document is property of Northeastern University.  Unauthorized distribution of the content is not~permitted.}  \\




\begin{enumerate}


\item 
Consider the following proposition, which we could label as $P(n)$. \\
 If $5n+3$ is even, then $n$ is odd. 

Prove by strong induction that for every positive integer $n$, proposition $P(n)$ is true. 

%Don't convert this to contrapositive form before using induction.   The purpose is to practice induction, as is.  \\

{\color{blue}
 $Proof:$\\
By strong induction, assume that $P(k)$ is true for $1\leqslant k \leqslant n-1$\\
The base case of $k=1$ holds as $P(1)=8$ which is even, and $1$ is odd.\\
As does the base case of $P(2)=13$, which is odd, and $2$ is even.\\
$5(n-2)+3=(5n+3)-10$\\
As $5n+3$ is even, as is 10, $5(n-2)+3$ is therefore even. \\
This proves the inductive hypothesis as $P(n-3)\implies P(n-2)$
$\blacksquare$

}



\newpage

\item Prove by strong induction that for every positive integer $n$ the following holds: 

$\sum\limits_{k=1}^n k(k{+}1) = \frac{n(n{+}1)(n{+}2)}{3}$

{\color{blue}
 $proof\colon$\\
 By strong induction, assume:\\
 $\sum\limits_{k=1}^{n-1} k(k{+}1) = \frac{n(n{+}1)(n-1)}{3}$\\
 The base case $n=1$ holds:\\
$\sum\limits_{k=1}^1 k(k{+}1) = 1(1+1)=2 = \frac{1(1{+}1)(1{+}2)}{3}$\\
 By subsituiting the last term of the summation:\\

 \sum\limits_{k=1}^n k(k{+}1) =\dfrac{n(n-1)(n+1)}{3}+n(n+1)$\\
 The inductive hypothesis is proved below:\\\\
 $\dfrac{n(n-1)(n+1)}{3}+n(n+1)=\dfrac{n(n+1)(n+2)}{3}\\
 n(n-1)(n+1)+3n(n+1)=n(n+1)(n+2)\\
 n^3 -n +3n^2 +3n = n^3 +3n +2n\\
 n^3+3n^2+2n=n^3+3n^2+2n\\
 \blacksquare$

}



\end{enumerate} 




\end{document}






















